\section{Kevin García: consulta de turnos mediante REST}

\subsection{Sistema y operación}
El sistema de Revisión Técnica Vehicular gestiona turnos, pagos e inspecciones. El personal utiliza una aplicación web conectada al backend. La operación analizada es la consulta del estado de un turno desde una aplicación para solicitantes.

\subsection{Arquitectura propuesta}
Se propone incorporar un servicio REST de consulta con una base de lectura independiente. El backend conservaría la base principal de turnos y pagos, y los cambios confirmados se propagarían mediante una cola de mensajes. Un consumidor actualizaría la base de lectura. La aplicación del solicitante consultaría el servicio REST mediante HTTP y recibiría una respuesta JSON.

Para publicar los cambios, el backend guardaría el turno y su evento en una misma transacción mediante una bandeja de salida. Un publicador enviaría los eventos registrados a la cola. El consumidor comprobaría el identificador y la versión del evento antes de actualizar la base de lectura, evitando procesar duplicados o reemplazar un estado reciente por otro anterior.

\subsection{Consulta REST}
La petición de ejemplo es \texttt{GET /api/turnos/123}. El servicio verificaría la identidad del solicitante y su permiso para consultar el turno. Después leería la base de consulta y devolvería \texttt{200 OK} con un JSON que incluiría el identificador, el estado y la fecha de actualización. El identificador 123 corresponde a un dato de prueba.

El servicio devolvería un error controlado cuando el turno no exista o el usuario no tenga acceso, sin exponer información de terceros. La aplicación mostraría el estado recibido junto con su fecha de actualización.

\subsection{Consistencia y manejo de fallos}
La propagación asíncrona permitiría atender consultas desde la base de lectura mientras se procesan los cambios de la base principal. Durante ese intervalo, el estado consultado podría corresponder a una versión anterior. La fecha de actualización permitiría identificar la antigüedad de la información.

Esta lectura se utilizaría para informar al solicitante. La confirmación de un pago y la autorización de una inspección se realizarían en el servicio responsable, consultando la información vigente.

Si el servicio REST no responde, la aplicación mostraría que la consulta está temporalmente no disponible y permitiría reintentar. Un resultado guardado previamente se mostraría con su fecha y con una indicación de que no se ha actualizado.

\subsection{Relación con CAP y PACELC}
Ante una partición que interrumpa la propagación de eventos, el servicio de consulta podría continuar respondiendo con los datos disponibles en su base de lectura, aceptando información temporalmente desactualizada. Las operaciones que requieran confirmar el estado vigente, como autorizar una inspección, se pospondrían cuando esa confirmación no sea posible.

Sin partición, PACELC permite analizar el equilibrio entre la latencia y la consistencia. La propuesta evita que cada consulta espere la propagación de todos los cambios y acepta un posible retraso del modelo de lectura. El tiempo de propagación y la latencia de consulta se evaluarían mediante mediciones.

\subsection{Criterio para elegir REST}
REST se propone para consultar un recurso identificado mediante HTTP y devolver su representación en JSON. La consulta no modifica el turno y permite integrar una aplicación cliente independiente. La cola de mensajes se utilizaría para propagar los cambios internos hacia la base de lectura.

SOAP se consideraría cuando una integración requiera operaciones descritas mediante WSDL y mensajes SOAP. Para la consulta planteada, se propone una interfaz HTTP con recursos identificados y respuestas JSON.

\subsection{Demostración de consulta con autorización}
La demostración utiliza el backend Spring Boot y PostgreSQL del RTV. Un cliente externo inicia sesión y envía \texttt{GET /api/consulta/turnos/82}. El backend obtiene el usuario de la sesión y comprueba en PostgreSQL la vinculación entre esa cuenta, el propietario y el turno. La cuenta no puede elegir el propietario mediante un parámetro de la petición.

La consulta del turno propio devuelve \texttt{200 OK} y los campos \texttt{turnoId}, \texttt{estado}, \texttt{fechaInicio} y \texttt{fechaFin}. Un turno ajeno o inexistente recibe el mismo \texttt{404}. Sin sesión se obtiene \texttt{401}; el intento de usar las rutas internas del personal devuelve \texttt{403} al solicitante. El acceso del personal autorizado se comprobó mediante su ruta interna y devolvió \texttt{200}.

Se prepararon una cuenta, un propietario y un turno ficticios para la prueba. Se conservaron las peticiones, las respuestas HTTP y las capturas, y pasaron diez pruebas automatizadas. Las fechas devueltas corresponden al turno. La cola de eventos, el servicio de lectura independiente y la propagación hacia su base corresponden a la arquitectura propuesta.

\subsection{Uso de inteligencia artificial}
Se utilizó ChatGPT Sol 6.1 en su forma Codex como apoyo para organizar y revisar la redacción de las notas y el diseño de la propuesta a partir del bosquejo de consulta de turnos.
